\section{Discussion}
\label{sec:discussion}

\subsection{What the measurements support}

The three backends agree on every instance of the exact suite and on every
instance of the campaigns they both solved: same canonical blocks, same LP
values, each with its own certificate produced by code that shares nothing with
the others. Together with the independent LP solver and the exhaustive rational
oracle, this supports the correctness of the implementation on the domain it
claims: generic DAGs, profits of either sign, strictly positive weights, one
capacity, cyclic digraphs by condensation.

On performance the picture is uneven, and it is given per family because a
single number would hide the shape of it. The parametric backend solves the whole
test set; Dinkelbach solves $40$ of $46$, the tuned \FS{} $36$ and the reference
\FSC{} $27$. The tuned \FS{} is not the fastest even where it finishes: its
median over solved instances, $0.142$\,s, is above the parametric backend's
$0.057$\,s, and it is taken over the smaller set it can finish.

Per family, the tuned \FS{} finishes everything on chains, diamonds, strongly
connected components, empty graphs, precedence forests, dense and sparse DAGs,
and the historical precedence-knapsack set; the parametric backend matches or
beats it on time in every one of those families --- a tie on the trivial ones,
a factor of $6$ on the historical instances and a factor of $29$ on the sparse
DAGs. It solves $3$ of $9$ layered
DAGs where the parametric backend solves all nine and Dinkelbach eight. On
disconnected graphs and on the MineLib-derived instances it solves nothing
within the limit, and neither does Dinkelbach; only the parametric backend does.
The claim that the \FS{} and Dinkelbach are comparable on the applied instances,
which an earlier reading of a handful of cases suggested, is not supported by
the frozen test set: on the historical precedence-knapsack family both finish
everything, and the medians are $0.087$\,s for Dinkelbach against $0.655$\,s.

One further correction belongs here rather than in a footnote. The ablation
shows two settings that the parameter search rejected and that are faster than
the frozen configuration on the paired comparison: the sink initial basis
($3.33$ against $2.83$) and transitive reduction enabled ($3.46$). The search
was run and frozen honestly, on training and validation data, but its choice is
not the best one available in the space it searched, and a second search with a
larger budget is the obvious next step.

\subsection{What changed, and why it matters more than the parameters}

The first measurements of this project had the \FS{} implementation one to three
orders of magnitude behind both baselines, with time-outs on medium instances.
The gap was not closed by tuning the pricing rule, which barely moves the
result, but by three structural corrections identified from the degeneracy
analysis of section~\ref{sec:robustness}: releasing the anti-cycling fallback
instead of letting it become the operating policy, restricting the ratio-test
scan to the arcs that can actually block, and stopping that scan at the first
zero-step blocker. The last two reproduce techniques already present in the
group's historical prototype, which had no anti-cycling rule at all and carried
an explicit feasibility shortcut for exactly this purpose.

The lesson recorded here is that in this problem class the entering rule is
close to irrelevant while the treatment of degeneracy dominates. Parameter
search over pricing schemes, block sizes and candidate limits produced
differences within a few per cent; the degeneracy-related settings produced
factors.

\subsection{What the measurements do not support}

They do not license any statement about the state of the art in maximum flow.
The flow backends use a Dinic implementation written for this project; they do
not use pseudoflow, highest-label push-relabel with gap and global relabelling,
or the amortised flow reuse of a proper parametric algorithm. A comparison
against a competitive flow code would very likely widen the gap on the families
where \FS{} already loses, and could also narrow the cases where it currently
wins.

They also do not support any claim of novelty. The audit of the literature has
not been carried out in this phase, by explicit decision, and
\file{docs/claims.md} keeps every novelty row in the \emph{to be verified}
state.

\subsection{What follows}

Two directions are supported by the measurements rather than by preference.
First, pricing and forest rebuilding are the dominant costs of a pivot, together
about four fifths of the solver's time on the hardest family, and both are
recomputed in full where only a small part of the graph changed: incremental
pricing and updates confined to the modified paths are the natural next step.
Second, and more important, the length of the pivot path is what actually
separates this method from the parametric backend, and neither the dual
formulation nor the leaving rules touched it.

A third direction is now ruled out rather than recommended. The dual simplex was
implemented and measured on a dedicated campaign precisely because the primal is
so degenerate, and it does remove the degeneracy it was predicted to remove; but
on the validation split it does not beat the primal, and the encouraging factors
observed on three hard instances did not survive the larger sample. That is the
reason the split exists.
